% Candidate statements and complete proofs for the DEG_L paper.
% Checked by /root/semireg_t7, GPT-6 Astra Pro; inherited effort not exposed.
% This is an input fragment, not a standalone article.
% Source locators are written explicitly so no guessed bibliography keys are used.

\subsection{The actual surface and its deformation problem}

\begin{proposition}[S-01: Allowed source covers do not preserve an embedding premise]
\label{prop:semireg-cover-domain}
Let $f:Y\to A$ be a surface morphism obtained from the four rational
Hodge maps in 032, Proposition~3.1. If $\pi:Y'\to Y$ is a connected finite
\'etale cover of degree $d>1$, the pulled-back maps still satisfy its
nonzero-period criterion, but $f\pi$ is not a closed immersion.
Moreover $(f\pi)_*[Y']=d f_*[Y]$.
\end{proposition}
\begin{proof}
Pullback preserves rational Hodge maps and multiplies the integral of
a top-degree class by $d$. Every geometric point of $Y$ has $d$ distinct
preimages with the same image under $f\pi$, so this morphism is not
injective on geometric points. The last assertion is the proper
pushforward formula for a finite map. The deeper principal covers
used here are explicitly permitted in 032, Corollary~6.7 and \S7.
\end{proof}

\begin{proposition}[S-02: The exact conormal gate]
\label{prop:semireg-dual-gate}
For a smooth closed immersion $i:Y\hookrightarrow A$ of a surface in
an abelian $g$-fold, $g\geq3$, write $N=N_{Y/A}$ and
$W=H^0(A,\Omega_A^1)$. Bloch semiregularity is equivalent to
surjectivity of
\[
 H^1(A,\mathcal O_A)\otimes\bigwedge^3W
 \longrightarrow H^1(Y,N^\vee\otimes K_Y).
\]
The map is restriction followed by the natural morphism
$\bigwedge^3 i^*\Omega_A^1\to N^\vee\otimes K_Y$.
\end{proposition}
\begin{proof}
Bloch's map is
$H^1(Y,N)\to H^{g-1}(A,\Omega_A^{g-3})$.
Serre duality identifies its dual source with $H^1(A,\Omega_A^3)$
and its dual target with $H^1(Y,N^\vee\otimes K_Y)$.
The conormal exact sequence induces the displayed exterior-power
map, since the quotient $\Omega_Y^1$ has rank two.
This is the defining dual map in Buchweitz--Flenner \S8.1;
Proposition~8.2 identifies it with Bloch's map.
Finally $\Omega_A^3=\mathcal O_A\otimes\bigwedge^3W$,
and finite-dimensional duality gives the equivalence.
\end{proof}

\begin{proposition}[S-03: A nonzero period does not force semiregularity]
\label{prop:semireg-period-counterexample}
There is a smooth surface in a CM abelian threefold with a nonzero
product of two holomorphic and two antiholomorphic one-forms for
which the Bloch semiregularity kernel has dimension at least eight.
\end{proposition}
\begin{proof}
Let $E$ be a CM elliptic curve, $B=E^2$, and
$L=\mathcal O_E(3[0])\boxtimes\mathcal O_E(3[0])$.
A general $C\in|L|$ is smooth and connected. Since $L^2=18$ and
$K_B=\mathcal O_B$, adjunction gives $g(C)=10$ and
$N_{C/B}=K_C$. Put $A=B\times E$ and $Y=C\times E$.
Then $N_{Y/A}=K_C\boxtimes\mathcal O_E$, so K\"unneth gives
$h^1(Y,N)=1+10=11$. The divisor semiregularity target
$H^2(A,\mathcal O_A)$ has dimension three.
The curve $C$ dominates the first elliptic factor: otherwise its
ample divisor class would be a fibre class. Hence a nonzero one-form
from that factor restricts to a nonzero form $\alpha$ on $C$.
For a nonzero one-form $\beta$ on the last $E$, the integral
$\int_Y(\alpha\wedge\beta)\wedge\overline{\alpha\wedge\beta}$
is nonzero, by the product formula and positivity of the two
one-dimensional integrals.
\end{proof}

\begin{proposition}[S-04: Rigid generating sources]
\label{prop:semireg-rigid-source}
Let $Y$ be a smooth projective surface with $H^1(Y,T_Y)=0$ and
$i:Y\hookrightarrow A$ a smooth closed immersion. Let $G\subset A$
be the abelian subvariety generated by a translate of $i(Y)$.
Then the ambient obstruction map satisfies
\[
\ker(\operatorname{ob}_i)=
\operatorname{Lie}(A)\otimes
\ker\bigl(H^1(A,\mathcal O_A)\to H^1(Y,\mathcal O_Y)\bigr).
\]
In particular it is injective when $G=A$. If in addition $i$ is
Bloch semiregular, this is also the infinitesimal Hodge-locus kernel
of the entire class $i_*[Y]$.
\end{proposition}
\begin{proof}
Restriction on $H^1(T_A)$ is $\operatorname{id}_{\operatorname{Lie}(A)}
\otimes i^*$, since $T_A$ is trivial. The normal sequence and
$H^1(T_Y)=0$ make $H^1(i^*T_A)\to H^1(N)$ injective.
This proves the kernel formula. The map $\operatorname{Alb}(Y)\to A$
has image $G$, so $i^*$ on $H^1(\mathcal O)$ is injective precisely
when $G=A$. Buchweitz--Flenner, Corollary~4.12 and Proposition~8.2,
identify
\[
 \sigma_i(\operatorname{ob}_i(\kappa))
 =\kappa\mathbin{\lrcorner}i_*[Y].
\]
If $\sigma_i$ is injective, the kernels of the two sides agree.
\end{proof}

The actual 032 source is a compact smooth complex two-ball quotient.
Calabi--Vesentini, \emph{Sur les vari\'et\'es complexes compactes
localement sym\'etriques}, Th\'eor\`eme~3, p.~315, gives
$H^1(Y,T_Y)=0$. The nonzero four-factor period does not establish
$G=A$. For example, separate injectivity on the two eigenspaces
of a quadratic action is insufficient: in
$A=E^n\times E^n$ with action $(a,\bar a)$, pullback to the diagonal
$E^n$ is injective on each eigenspace, but their images coincide.

\begin{proposition}[S-05: A quantitative obstruction for the actual morphism]
\label{prop:semireg-map-dimension}
Let $Y$ be a compact smooth complex two-ball quotient and
$f:Y\to A$ a morphism with two-dimensional image, where $\dim A=g$.
Let $\tau_f$ be the Buchweitz--Flenner morphism semiregularity map.
Then
\[
 \dim\ker\tau_f\ \geq\ gq(Y)-\binom{2g}{g-2}.
\]
\end{proposition}
\begin{proof}
In characteristic zero $df:T_Y\to f^*T_A=\mathcal O_Y^{\oplus g}$
is generically injective. Its kernel is a torsion-free sheaf of
generic rank zero, hence vanishes. With $N_f=\operatorname{coker}(df)$,
the derived dual of $L_{Y/A}$ is the two-term complex
$[T_Y\to f^*T_A]$ in degrees $0,1$, so
$T^2_{Y/A}(\mathcal O_Y)=H^1(Y,N_f)$.
Calabi--Vesentini and the cokernel sequence give
$H^1(\mathcal O_Y)^{\oplus g}\hookrightarrow H^1(N_f)$.
The target specified in Buchweitz--Flenner, Proposition~7.21,
Lemma~7.22 and the definition before Theorem~7.23, is
\[
 \bigoplus_{r=0}^{g-2}H^{r+2}(A,\Omega_A^r).
\]
Its dimension is
$\sum_{r=0}^{g-2}\binom gr\binom g{r+2}=\binom{2g}{g-2}$
by Vandermonde. Rank--nullity proves the assertion.
\end{proof}

\begin{corollary}[S-06: Failure at sufficiently deep allowed levels]
\label{cor:semireg-level-failure}
Fix a nonzero 032 surface construction $f:Y\to A$. There are
connected finite \'etale covers $\pi_j:Y_j\to Y$, preserving its
four pulled-back Hodge maps and nonzero period, for which
$\tau_{f\pi_j}$ is not injective. Such covers may be chosen among
a cofinal system of principal congruence levels.
\end{corollary}
\begin{proof}
The nonzero period implies $H^1(Y,\mathbf C)\ne0$.
Venkataramana, \emph{Virtual Betti numbers of compact locally
symmetric spaces}, Theorem~1, proves that finite-dimensional
direct-limit cohomology over finite covers equals the cohomology
of the compact dual. Here the compact dual is $\mathbf P^2$,
whose first cohomology is zero. Pullback under finite covers is
injective over $\mathbf Q$ by transfer, so the direct limit is
nonzero and therefore infinite-dimensional. Hence $q(Y_j)$ is
unbounded, and the preceding proposition applies with fixed $A$.
The congruence version is explicitly identified in the final Remark
on p.~3 of that paper. Its proof uses the same rational-conjugation
action, since rational conjugation and finite intersection preserve
congruence subgroups; principal subgroups are cofinal.
Persistence of the period follows from the first proposition.
\end{proof}

This is a statement about the deformation problem of a morphism.
The factorization
$\tau_f=\sigma_{Rf_*\mathcal O_Y}\circ\psi_f$
does not imply that a kernel of $\tau_f$ survives in the perfect
complex obstruction space: $\psi_f$ itself can have a kernel.
Nor does the corollary settle Bloch semiregularity of a particular
smooth reduced image.

\begin{proposition}[S-07: Descent of the cover obstruction space]
\label{prop:semireg-cover-invariants}
For a finite \'etale Galois cover $\pi:Y'\to Y$ with deck group $H$,
$N_{f\pi}=\pi^*N_f$ and
$H^1(Y',N_{f\pi})^H=H^1(Y,N_f)$.
\end{proposition}
\begin{proof}
Flat pullback identifies the tangent cokernel sequences.
Finite pushforward and the projection formula identify
$\pi_*\pi^*N_f=N_f\otimes\pi_*\mathcal O_{Y'}$.
In characteristic zero averaging over $H$ splits off precisely
the invariant summand $N_f$. Cohomology commutes with this direct
summand, giving the assertion.
\end{proof}

\subsection{Derived replacements and their remaining gate}

\begin{proposition}[S-09: An integral Chern-component projector]
\label{prop:semireg-k-projector}
Let $A$ be an abelian $g$-fold and $\xi\in K_0(A)$ with
$\operatorname{ch}_j(\xi)=0$ for $j<p$. For $m\geq2$ put
\[
 M=[m]^*,\qquad
 D=\prod_{j=p+1}^g(m^{2j}-m^{2p}),\qquad
 \eta=\prod_{j=p+1}^g(m^{2j}\operatorname{id}-M)\xi.
\]
Then $\eta\in K_0(A)$ and
$\operatorname{ch}(\eta)=D\operatorname{ch}_p(\xi)$ in rational
Betti cohomology.
\end{proposition}
\begin{proof}
The operator $M$ multiplies $H^{2j}(A,\mathbf Q)$ by $m^{2j}$.
Thus every higher Chern component is annihilated by its corresponding
factor, lower components were already zero, and the component of
index $p$ is multiplied by $D$. The polynomial has integer
coefficients. Integral K-classes are represented by perfect complexes
on the smooth projective variety, using finite locally free
resolutions and shifts to implement negative coefficients.
\end{proof}

Starting from an integral codimension-$p$ cycle of class $N_0t$,
one may use $\xi=\sum a_i[\mathcal O_{Z_i}]$; its leading Chern
character is the fundamental cycle. Alternatively, resolve each
$Z_i$ and use the proper pushforward of its structure sheaf;
Grothendieck--Riemann--Roch gives the same leading component.
Adding line-bundle classes gives the useful character
$Nt+\sum_a b_a e^{a\theta}$, which remains Hodge on the polarized
Route-2 family. This constructs a K-class, not a semiregular object.

\begin{proposition}[S-08: Multiplicity creates trace kernels]
\label{prop:semireg-multiplicity}
If $E$ is perfect and $\operatorname{Ext}^2(E,E)\ne0$, then
$E^{\oplus m}$ is not fully semiregular for any $m\geq2$.
\end{proposition}
\begin{proof}
Its degree-two self-Ext group is
$\operatorname{Ext}^2(E,E)\otimes\operatorname{End}(\mathbf C^m)$.
The Atiyah class is $\operatorname{at}(E)\otimes\operatorname{id}$;
therefore every semiregularity component is the corresponding
component for $E$ tensored with the matrix trace.
A nonzero degree-two Ext class tensored with a nonzero traceless
matrix lies in the kernel.
\end{proof}

\begin{proposition}[S-10: Full Fourier--Mukai transport]
\label{prop:semireg-full-fm}
For the fixed enhanced Fourier--Mukai equivalence in the supplied
032 companion \S5.2, the full trace maps satisfy
\[
 \sigma_{\Phi Q}\circ\Phi_{\rm Ext}
   =\Phi^H_{HH}\circ\sigma_Q
 \quad\hbox{on }\operatorname{Ext}^2_A(Q,Q).
\]
Thus full semiregularity is invariant under this equivalence.
\end{proposition}
\begin{proof}
Use the companion's enhanced adjunctions and notation.
Put
$b_Q(v)=\operatorname{Tr}_{Q,H_A[2]}(A_Q^H[2]\circ v)
\in HH_{-2}(A)$.
Equations (5.20), (5.24) identify its HKR image with $\sigma_Q(v)$.
For $\delta\in T^{-2}(A)$, trace-compatible adjunction gives
\[
 \langle b_Q(v),\delta\rangle
   =\operatorname{Tr}_{A,Q}(\delta_Q[2]\circ v).
\]
Indeed, if $\delta':H_A\to\omega_A[g-2]$ is the adjoint coefficient
map, move $\delta'[2]$ through the partial trace by (5.24);
evaluation then gives the right side. This is exactly (5.37)
with the displayed even shift.
The convolution proof of (5.33), with the same inverse adjunction
data and shifted by two, gives
$(\Psi^{-2}\delta)_{\Phi Q}=\theta_Q[-2]\Phi(\delta_Q)$.
By the trace preservation in (5.31),
\[
 \langle b_{\Phi Q}(\Phi_{\rm Ext}v),\Psi^{-2}\delta\rangle
  =\operatorname{Tr}_{A,Q}(\delta_Q[2]\circ v)
  =\langle b_Q(v),\delta\rangle.
\]
Definition (5.34) and nondegeneracy of the Serre pairing imply
$b_{\Phi Q}(\Phi_{\rm Ext}v)=\Phi_*^{HH}b_Q(v)$.
Apply (5.35) and the fixed HKR maps. The two vertical maps
are isomorphisms, proving the assertion about kernels.
This proves a refinement; Proposition~5.2 itself only asserted
transport on the obstruction-evaluation image.
\end{proof}

\begin{proposition}[S-11: An ample complete intersection need not be semiregular]
\label{prop:semireg-ci-counterexample}
Let $A$ be an abelian fourfold, $L$ very ample, and
$Y=D_1\cap D_2$ a smooth transverse intersection of two members
of $|L|$. Then $h^1(N_{Y/A})=24$, whereas
$h^3(A,\Omega_A^1)=16$. Thus its Bloch map has a kernel of
dimension at least eight.
\end{proposition}
\begin{proof}
The twisted Koszul sequence is
$0\to L^{-1}\to\mathcal O_A^{\oplus2}\to\mathcal I_Y\otimes L\to0$.
Kodaira vanishing and $K_A=\mathcal O_A$ give $H^j(A,L)=0$ for
$j>0$; Serre duality gives $H^j(A,L^{-1})=0$ for $j<4$.
Thus
$H^2(A,\mathcal I_Y\otimes L)=H^2(A,\mathcal O_A)^{\oplus2}$.
The restriction sequence gives
$H^1(Y,L|_Y)=H^2(A,\mathcal I_Y\otimes L)$, of dimension twelve.
Transversality identifies $N_{Y/A}=(L|_Y)^{\oplus2}$.
The Bloch target has dimension
$\binom43\binom41=16$, so rank--nullity finishes the proof.
\end{proof}

\begin{proposition}[S-12: A pure Weil character obstruction for coherent sheaves]
\label{prop:semireg-pure-weil-no-go}
Let $A_0$ be a projective member of Voisin's four-dimensional
marked $\mathbf Q(i)$-Weil torus family, and let $t\ne0$ be a
rational Weil class. No fully Buchweitz--Flenner-semiregular
coherent sheaf $F$ on $A_0$ satisfies
\[
 \operatorname{ch}_1(F)=\operatorname{ch}_3(F)=0,\qquad
 \operatorname{ch}_2(F)=t.
\]
The rank and top-degree character are unrestricted.
\end{proposition}
\begin{proof}
The unpolarized Weil torus parameter space is smooth near $A_0$;
Voisin, \emph{A counterexample to the Hodge conjecture extended
to K\"ahler varieties}, \S3, equations (3.1)--(3.8), describes it.
Lemma~3.1 shows that the rational Weil space remains of Hodge
type $(2,2)$. Hence the components with indices $I=\{1,2,3\}$
in the assertion remain horizontal Hodge classes. In dimension
four these are precisely the components tested by full
semiregularity.
Buchweitz--Flenner Theorem~5.1 and its proof on p.~179 would
therefore extend $F$ as a coherent flat sheaf on a neighborhood.
Voisin Proposition~3.2 and Proposition~2.1 show that very general
fibres in that neighborhood have $\operatorname{NS}=0$ and
$c_2(G)=0$ for every coherent analytic sheaf $G$.
On such a fibre $c_1(F)=c_2(F)=0$, whence
$\operatorname{ch}_2(F)=0$, contrary to the nonzero horizontal
class $t$.
\end{proof}

This last proposition is a coherent-sheaf statement for the
specified fourfold family. Extension to arbitrary perfect complexes
would require an analytic effectivity theorem for their formal
deformations; algebraic/formal Pridham alone is not that theorem.
Allowing a nonzero polarization term in the Chern character
avoids this particular contradiction.

\begin{proposition}[S-13: One exponential polarization term does not repair the coherent gate]
\label{prop:semireg-single-exponential-no-go}
In the fourfold setting of S-12, let $L$ be a polarization whose
Rosati involution restricts to conjugation on $\mathbf Q(i)$, and
write $\theta=c_1(L)$. If $a\in\mathbf Z$, $\lambda\in\mathbf Q^\times$,
and $t\ne0$ is a rational Weil class, no fully semiregular coherent
sheaf $F$ has
\[
 \operatorname{ch}(F)=r e^{a\theta}+\lambda t+c[\mathrm{pt}]
\]
for any $r,c\in\mathbf Q$.
\end{proposition}
\begin{proof}
Write $H^1(A_0,\mathbf C)=V_\sigma\oplus V_{\bar\sigma}$.
The Rosati condition places $\theta$ in
$V_\sigma\otimes V_{\bar\sigma}$, whereas the complexified Weil
space is
$\bigwedge^4V_\sigma\oplus\bigwedge^4V_{\bar\sigma}$.
Both eigenspaces have dimension four. Consequently
$\theta\smile t=0$, and hence $e^{-a\theta}t=t$.
Tensoring by $L^{-a}$ is a Fourier--Mukai autoequivalence.
The functoriality in Perry Remarks~2.3 and~2.5, p.~5, preserves
full semiregularity under this equivalence, as in S-10.
It preserves coherence as well.
For $F'=F\otimes L^{-a}$ the Chern-character formula gives
\[
 \operatorname{ch}(F')=r+\lambda t+c[\mathrm{pt}].
\]
Thus $\operatorname{ch}_1(F')=\operatorname{ch}_3(F')=0$ and
$\operatorname{ch}_2(F')=\lambda t\ne0$, contradicting S-12.
\end{proof}

The integer condition on $a$ ensures that this normalization is an
actual line-bundle twist. The conclusion does not exclude all
polarization polynomials: it excludes the displayed single
exponential, including an arbitrary additional top-degree term.

\begin{proposition}[S-14: An equivariant perfect-complex gate implies bounded division]
\label{prop:semireg-perry-bounded-division}
Let $\pi:A\to S$ be an abelian scheme over a smooth irreducible
complex algebraic variety of finite type, with a relatively very
ample line bundle $H$ and $\theta=c_1(H)$. Let
$t\in H^0(S,R^{2p}\pi_*\mathbf Q(p))$ be a flat section whose values
are Hodge classes, where $0\leq p\leq\dim(A/S)$.
Suppose that for one $s_0\in S(\mathbf C)$ there exist:
\begin{enumerate}
\item a finite subgroup $G\subset A_{s_0}\times A_{s_0}^{\vee}$
with a specified coherent action on $\operatorname{Perf}(A_{s_0})$
by translations and tensoring by line bundles;
\item a weakly $G$-semiregular perfect complex $E_0$ with
$\operatorname{Ext}^{i}(E_0,E_0)=0$ for all $i<0$;
\item an algebraic class $B_0\in H^2(A_{s_0},\mathbf Q(1))$,
a nonzero $\lambda\in\mathbf Q$, and $P\in\mathbf Q[x]$ such that
\[
 e^{B_0}\operatorname{ch}(E_0)=\lambda t_{s_0}+P(\theta_{s_0})
\]
in the full even rational cohomology.
\end{enumerate}
Then there is an integer $C\geq1$ such that, for every
$s\in S(\mathbf C)$, there are an integer $1\leq N_s\leq C$
and effective integral codimension-$p$ cycles $Z_s^+,Z_s^-$,
each of $H$-degree at most $C$, with
\[
 N_s t_s=\operatorname{cl}(Z_s^+)-\operatorname{cl}(Z_s^-).
\]
In particular, there are single integers $N,D\geq1$ for the entire
family such that $Nt_s$ admits such representatives of degrees
at most $D$ on every fibre.
A sufficient form of the second hypothesis is a specified
$G$-equivariant structure on $E_0$ for which the ordinary full
semiregularity map is injective on
$\operatorname{Ext}^2(E_0,E_0)^G$, together with the same vanishing
of negative self-Ext groups.
\end{proposition}
\begin{proof}
The class
$w=\lambda t+P(\theta)$ is a global flat section of the full even
cohomology local system, and every component has the required
Hodge type. It extends
$w_0=e^{B_0}\operatorname{ch}(E_0)$.
Apply Alexander Perry, \emph{The semiregularity theorem for
equivariant noncommutative varieties},
arXiv:2604.00511v2 (24 June 2026), Theorem~1.2, p.~3.
Its second conclusion states that $w_s$ is algebraic for every
$s\in S(\mathbf C)$. Since $P(\theta_s)$ is algebraic and
$\lambda\ne0$, each $t_s$ is algebraic.
DIV1(2), applied to the fixed smooth projective family and $H$,
now supplies a single bound $C$ for the signed integral complexity.
This is exactly the displayed conclusion.
For a common multiplier take $N=\operatorname{lcm}(1,\ldots,C)$
and multiply the cycles on fibre $s$ by $N/N_s$.
Their degrees are at most $D=NC$.
Finally, Perry Definition~2.6 and Lemma~2.7, pp.~5--6, show that
the stated invariant-injectivity criterion implies weak
$G$-semiregularity over $\mathbf C$.
\end{proof}

The gate requires the entire displayed Chern character to remain
Hodge, a coherent categorical action, and vanishing of negative
self-Ext groups. A virtual K-class alone supplies none of these
deformation properties. The theorem supplies algebraicity on
every fibre; no properness assertion for a moduli space of
perfect complexes is used. The bound obtained from DIV1 is
existential. Here $G$ is a subgroup of the target's translation
and degree-zero line-bundle group, not an arbitrary deck group
of the source surface.

\begin{proposition}[S-15: Canonical source-cover equivariance retains a full kernel]
\label{prop:semireg-cover-equivariance}
Let $f:Y\to A$ be a proper morphism from a smooth projective
variety to an abelian variety, and put $E=Rf_*\mathcal O_Y$.
\begin{enumerate}
\item If $\rho:Y'\to Y$ is a finite \'etale Galois cover with
group $K$, the canonical $K$-equivariant lift of
$E'=R(f\rho)_*\mathcal O_{Y'}$ for the trivial action of $K$
on $\operatorname{Perf}(A)$ can be semiregular only if $E$
is fully semiregular.
\item Suppose $\rho$ is obtained by base change from an isogeny
$q:A'\to A$. Set
$G=\ker(q^*:A^\vee\to(A')^\vee)$, acting on
$\operatorname{Perf}(A)$ by the coherently chosen tensor
line bundles in $q_*\mathcal O_{A'}$.
The canonical $G$-equivariant lift of
$q_*q^*E=R(f\rho)_*\mathcal O_{Y'}$ can be semiregular only if
$E$ is fully semiregular.
\end{enumerate}
These assertions concern the canonical lifts supplied by the
covers; they do not exclude every possible new equivariant
structure on the underlying objects.
\end{proposition}
\begin{proof}
For (1), finite pushforward and characteristic-zero averaging give
\[
 (E')^K
 =Rf_*(\rho_*\mathcal O_{Y'})^K
 =Rf_*\mathcal O_Y=E.
\]
Because the categorical action is trivial,
$\operatorname{Perf}(A)^K$ is the finite product of copies of
$\operatorname{Perf}(A)$ indexed by irreducible representations
of $K$. The factor belonging to the trivial representation of
the canonical lift is $E$. Self-Ext groups, Hochschild homology,
and the semiregularity map decompose by these factors.
Semiregularity of the lift therefore implies that of $E$.

For (2), let $R=q_*\mathcal O_{A'}$. The character decomposition
of the finite abelian torsor gives
$R=\bigoplus_{P\in G}P$, with multiplication isomorphisms
$P\otimes Q\cong PQ$. These provide the indicated coherent
tensor action. A $G$-equivariant object is equivalently an
$R$-module object; since $q$ is finite \'etale this identifies
$\operatorname{Perf}(A)^G$ with $\operatorname{Perf}(A')$.
Under this equivalence the canonical lift of $q_*q^*E$
corresponds to $q^*E$.
Flat base change for the square defining $Y'$ identifies the
underlying object with the stated source-cover pushforward.

The pullback map
\[
 q^*:\operatorname{Ext}^2_A(E,E)
 \longrightarrow
 \operatorname{Ext}^2_{A'}(q^*E,q^*E)
\]
is injective: by adjunction its target is
$\operatorname{Ext}^2_A(E,E\otimes R)$, and the unit
$\mathcal O_A\to R$ has the retraction
$(\deg q)^{-1}\operatorname{Tr}_q$.
Semiregularity is natural for the enhanced functor $q^*$,
by Perry Remark~2.3, p.~5. Thus pullback carries
$\ker\sigma_E$ injectively into $\ker\sigma_{q^*E}$.
Semiregularity of the canonical lift would make the latter
kernel zero, proving the assertion.
\end{proof}

For an ordinary source deck transformation $k$ one has
$f\rho k=f\rho$. Consequently an induced translation of the
target satisfying this equivariance must be the identity.
An ordinary deck action therefore supplies internal
automorphisms of the pushed-forward object; it does not by
itself specify the nontrivial target action required in S-14.
Part (2) treats a special case where target tensor symmetries
really do arise, and shows why their canonical lift still
does not remove a pre-existing full semiregularity kernel.

There is also an independent primary-source confirmation of S-10:
Perry Remarks~2.3 and~2.5, p.~5, give functoriality of the
semiregularity morphism on the entire mapping complex and identify
its degree-two part with Buchweitz--Flenner under HKR.
An equivalence induces isomorphisms on both mapping complexes and
Hochschild homology, so the full kernel is invariant. This argument
does not depend on a raw versus Todd-modified HKR convention;
for abelian varieties the Todd class is one.

\begin{proposition}[S-16: A polarization exponential realizes a K-class by a vector bundle]
\label{prop:semireg-kclass-bundle-realization}
Let $X$ be a smooth projective variety, let $L$ be ample, and let
$\xi\in K_0(X)$. There exist integers $m,r\geq1$ and a vector
bundle $V$ such that
\[
 [V]=\xi+r[L^{-m}]\quad\hbox{in }K_0(X).
\]
In particular, if $\operatorname{ch}(\xi)=\lambda t$, then
$\operatorname{ch}(V)=r e^{-m c_1(L)}+\lambda t$.
In the setting of S-13, with $\lambda t\ne0$, every vector bundle
with this character fails full semiregularity.
\end{proposition}
\begin{proof}
On a smooth projective variety a K-class is a difference
$\xi=[U]-[W]$ of vector bundles: take finite locally free
resolutions and group their even and odd terms.
For sufficiently large $m$, the vector bundle $W\otimes L^m$
is generated by finitely many global sections.
The precise ampleness criterion is Stacks Project,
Proposition~28.27.13(7)--(8), tag 01Q3.
Choose a surjection
\[
 0\longrightarrow K\longrightarrow\mathcal O_X^{\oplus r}
   \longrightarrow W\otimes L^m\longrightarrow0 .
\]
One can add a redundant zero section to ensure $r\geq1$.
The kernel $K$ is locally free, since the surjection onto a
vector bundle splits locally. After tensoring by $L^{-m}$,
the K-group relation is
$[K\otimes L^{-m}]=r[L^{-m}]-[W]$.
Thus $V=U\oplus(K\otimes L^{-m})$ has the required class.
The Chern-character formula follows by additivity, and the
last assertion is S-13 with $a=-m$ and $c=0$.
\end{proof}

This gives an actual coherent vector-bundle representative
after one exponential addition. It supplies neither full nor
equivariant semiregularity; in the indicated fourfold case the
full version is ruled out by S-13.

\begin{proposition}[S-17: Tensor--determinant untwisting after a positive-rank deformation]
\label{prop:semireg-tensor-determinant}
Let $\pi:\mathcal X\to X$ be a $\mu_M$-gerbe over a connected
smooth complex scheme, and let $F$ be a weight-one perfect complex
of positive virtual rank $r$. Then
\[
 F^{\otimes^{\mathbf L}r}\otimes\det(F)^{-1}
 \simeq\pi^*K
\]
for an ordinary perfect complex $K$ on $X$.

In the family setting and under all the hypotheses of S-14,
suppose $\dim A_s=2n$, $n\ge2$, that $t$ is a nonzero relative
Weil class of codimension $n$, and that
\[
 w_0=e^{B_0}\operatorname{ch}(E_0)
     =P(\theta_0)+\lambda t_0,
 \qquad P(0)=r>0,\quad\lambda\ne0,\quad\theta_0t_0=0.
\]
Put $p_1=[x]P(x)$. On the local deformation neighborhood supplied
by Perry, the descended complex has character, through
codimension $n$,
\[
 \operatorname{ch}(K_s)
 \equiv e^{-p_1\theta_s}P(\theta_s)^r+r^r\lambda t_s.
 \tag{S17.1}
\]
There are fixed integers $R,m\ge1$ and $a,k\ge0$ from a finite
bundle presentation and generation twists for which the following
numbers give a uniform signed-cycle bound on the entire base.
Set
\[
 Q(x)=e^{kx}\bigl(e^{-p_1x}P(x)^r+R e^{-mx}+a\bigr),
 \qquad\mu=r^r\lambda,
\]
\[
 C_n=[u^n]\exp\!\left(\sum_{j=1}^n
     (-1)^{j-1}(j-1)![x^j]Q(x)u^j\right),
 \qquad\beta=(-1)^{n-1}(n-1)!\mu .
\]
If $d\ge1$ clears the denominators of $C_n$ and $\beta$, and
$v=\int_{A_s}\theta_s^{2n}$, then
\[
 N=|d\beta|>0,\qquad D=dC_nv>0
 \tag{S17.2}
\]
are fixed integers such that $Nt_s$ is the difference of two
effective integral cycles, each of polarization degree at most
$D$, on every fibre. A negative nonzero virtual rank is reduced
to this case by shifting the initial and deformed complexes.
\end{proposition}
\begin{proof}
The determinant functor for perfect complexes exists and commutes
with base change by Knudsen--Mumford, \emph{The projectivity of the
moduli space of stable curves I}, Theorem~2, p.~42, with the
extended determinant definition on pp.~40--42. Apply it on an
\'etale atlas of the gerbe. Functoriality and base-change
compatibility glue the local determinants to a line bundle
$\det(F)$ on $\mathcal X$. Locally a finite locally free
presentation gives
\[
 \det(F)=\bigotimes_i\det(F^i)^{(-1)^i}.
\]
The inertia element $\zeta\in\mu_M$ acts on this line by
$\zeta^{\sum_i(-1)^i\operatorname{rk}F^i}=\zeta^r$.
Consequently $F^{\otimes^{\mathbf L}r}\otimes\det(F)^{-1}$
has weight zero. Weight-zero quasi-coherent complexes on a
$\mu_M$-gerbe are precisely pullbacks from the base: after an
\'etale trivialization this is the weight-zero summand of the
representation category of $\mu_M$, and the identifications
descend. Perfectness is local for this topology, proving the
first assertion.

Use the weight-one deformation $F$ on the gerbe in Perry's proof
of Theorem~1.2, \S6.2, pp.~26--28. At the central fibre its
normalization has character $w_0$ and first component
$w_{0,1}=p_1\theta_0$. Directly, cancellation of the $B$-field
in the tensor--determinant expression gives
\[
 \operatorname{ch}(K_0)
 =\operatorname{ch}(E_0)^r
       e^{-c_1(E_0)}
 =w_0^r e^{-w_{0,1}} .
\]
This also follows from the central root description of the gerbe:
the root-line factors cancel between the $r$ tensor factors and
the determinant. Since $n\ge2$, $w_{0,1}=p_1\theta_0$ contains
no Weil term. Through codimension $n$, the only occurrence of
$t_0$ in $w_0^r$ is
$rP(0)^{r-1}\lambda t_0=r^r\lambda t_0$.
All positive powers of $\theta_0$ annihilate $t_0$.
This proves (S17.1) at the central fibre; flatness of Chern classes
proves it over the connected deformation neighborhood.
Terms in higher codimension, including a possible $t_0^2$ term,
are not used in this formula or in the computation of $c_n$.

Shrink the neighborhood to a smooth affine base $U$.
Its total space $A_U$ is smooth and quasi-projective, and its
relatively very ample bundle $H$ is ample over this affine base.
Stacks Project Lemma~36.37.2, tag 0F8F, represents a perfect
complex by a bounded complex of vector bundles on a scheme
with the resolution property. Thus write
$[K]=[V^+]-[V^-]$. For sufficiently large $m$, choose a finite
surjection $\mathcal O^{\oplus R}\to V^-\otimes H^m$.
Its kernel $J$ is locally free, and
\[
 [V^+\oplus(J\otimes H^{-m})]=[K]+R[H^{-m}].
\]
Adding redundant sections ensures $R\ge1$.
This is the relative version of S-16, using the finite global
generation criterion of Stacks Proposition~28.27.13(7)--(8),
tag 01Q3. Add $a$ trivial summands to make the rank at least $n$,
and twist by a sufficiently large $H^k$ to obtain a globally
generated vector bundle $V$. Its character through codimension
$n$ is $Q(\theta_s)+\mu t_s$.

The identity
\[
 c_u(V)=\exp\!\left(\sum_{j\ge1}
             (-1)^{j-1}(j-1)!\operatorname{ch}_j(V)u^j\right)
\]
gives $c_n(V_s)=C_n\theta_s^n+\beta t_s$.
If $q=\operatorname{rk}V$, the degeneracy cycle of $q-n+1$
general sections is effective of codimension $n$, or empty,
and represents $c_n(V_s)$; see Fulton, \emph{Intersection
Theory}, Theorem~14.4 and Examples~14.4.2--3, pp.~254--256.
Call it $Z_s$. Its Weil component $\beta t_s$ is nonzero,
so the cycle is nonempty. Since $\theta_st_s=0$, its degree is
$C_nv$, which is positive. Let $W_s$ be a complete intersection
of $n$ general members of $|H_s|$; it represents $\theta_s^n$
and has degree $v$. Then
\[
 d\beta t_s=\operatorname{cl}(dZ_s)
                   -\operatorname{cl}(dC_nW_s).
\]
Both cycles are integral and effective of degree $D$.
If $\beta<0$, interchange them. This proves (S17.2) over a
nonempty \'etale image, which is a dense Zariski open subset
of the irreducible base. The fixed-$N$, fixed-$D$ locus of such
signed representatives is Zariski closed by the proper
parameter-space argument in DIV1. It therefore equals the
entire base, preserving these same constants.

Finally, a shift changes the signs of the virtual rank and
Chern character, preserves the self-Ext groups and
semiregularity, and shifts the existing deformation.
This reduces negative rank to positive rank without adding
a new existence hypothesis.
\end{proof}

This construction is used only after the successful twisted
deformation exists. No semiregularity of $K$ or of $V$ is claimed.
The Adams construction DER4 removes the nonzero-rank restriction;
S-17 is a separate, simpler tensor--determinant alternative.

\begin{proposition}[S-18: Canonical induction preserves a full semiregularity kernel]
\label{prop:semireg-induction-kernel}
Let $\mathcal C$ be a smooth proper complex linear enhanced
category with a coherent action of a finite group $G$ by exact
autoequivalences, and assume its equivariant category
$\mathcal C^G$ is smooth proper. Let
\[
 \operatorname{Ind}:\mathcal C\longrightarrow\mathcal C^G
\]
be induction, whose underlying object is
$\bigoplus_{g\in G}gE$, with the canonical permutation
equivariant structure. Induction is injective on every morphism
space. In particular it induces an injection
\[
 \ker\sigma_E\hookrightarrow
       \ker\sigma_{\operatorname{Ind}(E)}.
\]
Thus a nonzero full semiregularity kernel cannot be removed by
canonical induction, including for actions combining
translations and tensoring by degree-zero line bundles on an
abelian variety.
\end{proposition}
\begin{proof}
For $u:E\to E'[j]$, the underlying map of
$\operatorname{Ind}(u)$ is the block diagonal map
$\bigoplus_{g\in G}g(u)$.
If it is zero, compose with the inclusion of the identity
summand of the source and the projection to the identity
summand of the target. The coherence isomorphism for the
identity element identifies this composite with $u$, so $u=0$.
In particular, if $\operatorname{Ind}(u)$ is zero in the
equivariant category, forgetting equivariance gives zero and
the same argument applies. This proves injectivity in every
degree, without requiring the forgetful functor to be full.

Induction is an exact enhanced functor, so naturality of the
full semiregularity morphism gives
\[
 \sigma_{\operatorname{Ind}(E)}(\operatorname{Ind}(u))
       =\operatorname{Ind}_*\sigma_E(u).
\]
This is Perry Remark~2.3, p.~5; Remark~2.5 identifies the
geometric map with Buchweitz--Flenner under HKR.
Consequently a nonzero element of $\ker\sigma_E$ has a
nonzero image lying in the displayed target kernel.
\end{proof}

S-18 rules out semiregularity of the canonical induced lift.
It does not rule out a different equivariant structure on the
same underlying object: weak equivariant semiregularity is
an existence condition over possible lifts, as specified in
Perry Definition~2.6, pp.~5--6.

\begin{proposition}[S-19: The isolated-point genus-four secant seed fails at every parameter]
\label{prop:semireg-markman-genus-four-points}
Consider the first genus-four construction of Markman
Example~8.2.3 in the supplied 31 May 2025 source. Thus $X$ is
a genus-four Jacobian with principal polarization $\theta$,
$d\ge1$, and $G\subset X$ is a cyclic translation group of
order $\ell=d+1$. The $G$-invariant reduced subscheme $Z$ is
the union of $\ell$ translates of $W_2$ in the stated general
position and $2d\ell$ isolated reduced points away from them.
Let $q:X\to Y=X/G$, let $\bar Z=Z/G$, and set
$E=\mathcal I_{\bar Z}$. Then
\[
 \chi(E,E)=8d,\qquad
 \operatorname{ext}^1(E,E)\ge8d+4,\qquad
 \operatorname{ext}^2(E,E)\ge24d+6>28.
 \tag{S19.1}
\]
In particular $E$ is not fully semiregular. The original ideal
$\mathcal I_Z$ is not weakly $G$-semiregular for this translation
action, for any choice of linearization. The same holds for
$\mathcal I_Z(\theta)$ under the action conjugated by tensoring
with $\mathcal O_X(\theta)$.
\end{proposition}
\begin{proof}
The quotient is finite \'etale of degree $\ell$. Its isolated
point set has $2d$ reduced points, and
$q^*\mathcal I_{\bar Z}=\mathcal I_Z$.
Markman's character calculation in Example~8.2.3 gives
\[
 w:=\operatorname{ch}(\mathcal I_Z(\theta))
   =1+\theta-\frac d2\theta^2-\frac d6\theta^3
          +\frac{d^2}{24}\theta^4 .
\]
Since $\int_X\theta^4=24$, direct multiplication yields
$\int_X w^\vee w=8d(d+1)$. Tensoring by the theta line bundle
does not change this Euler pairing. Riemann--Roch and the
\'etale degree formula therefore give $\chi(E,E)=8d$.

We verify the asserted tangent-space lower bound. On a smooth
variety, the endomorphism sheaf of an ideal of codimension at
least two is $\mathcal O$: an endomorphism is multiplication
by a regular function away from that locus, and normality
extends the function. The unit and perfect trace give a split
inclusion
\[
 \mathcal O_Y\longrightarrow R\mathcal Hom(E,E)
       \longrightarrow\mathcal O_Y
\]
whose composite is the virtual rank, namely one. Its complement
has no degree-zero cohomology and no negative cohomology.
Consequently
\[
 \operatorname{Ext}^1(E,E)
 \cong H^1(Y,\mathcal O_Y)
       \oplus H^0(Y,\mathcal Ext^1(E,E)).
\]
At each isolated reduced point the ideal stalk is the maximal
ideal $\mathfrak m$ of a regular local ring $R$ of dimension
four. Apply $\operatorname{Hom}_R(\mathfrak m,-)$ to
$0\to\mathfrak m\to R\to\mathbf C\to0$.
Here $\operatorname{Hom}_R(\mathfrak m,R)=R$ maps to zero in
$\operatorname{Hom}_R(\mathfrak m,\mathbf C)$, and
\[
 \operatorname{Ext}^1_R(\mathfrak m,R)
   =\operatorname{Ext}^2_R(\mathbf C,R)=0.
\]
Thus
$\operatorname{Ext}^1_R(\mathfrak m,\mathfrak m)
\cong(\mathfrak m/\mathfrak m^2)^\vee$ has dimension four.
The isolated supports contribute independent skyscraper
summands. Since $h^1(\mathcal O_Y)=4$, this proves
$\operatorname{ext}^1(E,E)\ge4+8d$.

The sheaf $E$ is simple. Serre duality on the abelian fourfold
gives $\operatorname{ext}^4=\operatorname{ext}^0=1$ and
$\operatorname{ext}^3=\operatorname{ext}^1$; there are no
self-Ext groups outside degrees zero through four. Hence
\[
 \operatorname{ext}^2(E,E)
 =\chi(E,E)-2+2\operatorname{ext}^1(E,E)
 \ge24d+6.
\]
The full semiregularity target on an abelian fourfold has
dimension
\[
 h^2(\mathcal O_Y)+h^3(\Omega_Y^1)+h^4(\Omega_Y^2)
 =6+16+6=28,
\]
proving failure of injectivity.

Equivariant descent identifies
$\operatorname{Perf}(X)^G$ with $\operatorname{Perf}(Y)$
and the canonical lift of $\mathcal I_Z$ with $E$.
Moreover $\operatorname{Aut}(\mathcal I_Z)=\mathbf C^*$,
so any two $G$-linearizations differ by a character of $G$.
Their descents differ by tensoring with the associated
degree-zero torsion line bundle on $Y$. Such a tensor
equivalence preserves the full kernel by S-10.
Thus no linearization is semiregular in the invariant category.
Conjugating this action by the theta tensor equivalence gives
the last assertion; the conjugated action consists of target
translations and tensoring by degree-zero line bundles.
\end{proof}

\begin{proposition}[S-20: The fourfold adjunction test excludes the partial-normalization seed]
\label{prop:semireg-fourfold-partial-normalization}
On an abelian fourfold $Y$, let $E$ be perfect and let
$c_E^2:HT^2(Y)\to H\Omega^{-2}(Y)$ be contraction with
$\operatorname{ch}(E)$, in the conventions of the supplied
032 companion, \S5.1. If $E$ is fully semiregular, then
\[
 \operatorname{ext}^2(E,E)=\operatorname{rank}c_E^2.
 \tag{S20.1}
\]
For the cyclic partial-normalization construction in Markman
Example~8.2.3, let $\bar E$ be the descended complex
\[
 [\mathcal O_Y\longrightarrow
       \bar\nu_*\mathcal O_{\widetilde Z/G}]
\]
in degrees zero and one, where four of the six transverse
meeting points for each pair of translated $W_2$ surfaces
are normalized, and the collection of choices is $G$-invariant.
Then, for every $d\ge1$,
\[
 \operatorname{rank}c_{\bar E}^2=12,
 \qquad
 \operatorname{ext}^2(\bar E,\bar E)\ge8d+6>12.
 \tag{S20.2}
\]
Thus this complex is not fully semiregular, and its pullback
is not weakly $G$-semiregular for the specified free translation
action, with any linearization. In fact
\[
 \dim\ker\sigma_{\bar E}\ge4d-3.
 \tag{S20.3}
\]
\end{proposition}
\begin{proof}
First we check the adjunction in (S20.1). Fix a nonzero volume
form on $Y$. Serre duality gives a perfect pairing on
$\operatorname{Ext}^2(E,E)$, while contraction and integration
give the perfect pairing
\[
 HT^2(Y)\otimes H\Omega^{-2}(Y)\longrightarrow\mathbf C.
\]
Indeed its $p$-th block pairs
$H^{2-p}(\wedge^pT_Y)$ with $H^{p+2}(\Omega_Y^p)$,
for $0\le p\le2$.
The evaluation identity (5.21) in the supplied 032 companion
is $\operatorname{ob}_E(\xi)=\xi\mathbin{\lrcorner}
\exp(a_E)$. Its equations (5.20) and (5.24), or directly the
definition of perfect trace with coefficients, then give
\[
 \int_Y\operatorname{Tr}
       \bigl(\operatorname{ob}_E(\xi)\circ u\bigr)
   =\langle\xi,\sigma_E(u)\rangle,
 \qquad u\in\operatorname{Ext}^2(E,E).
\]
The volume form is implicit in the integral. Both morphisms
have degree two, so no odd exchange sign is introduced.
The factorial in each block is the same $1/p!$ in evaluation
and in the trace map. Thus $\sigma_E$ is the transpose of
$\operatorname{ob}_E:HT^2(Y)\to\operatorname{Ext}^2(E,E)$.
In particular injectivity of $\sigma_E$ is equivalent to
surjectivity of $\operatorname{ob}_E$.
The entire identity
$\sigma_E\operatorname{ob}_E=c_E^2$ is Proposition~5.1,
equation (5.9), of that companion; it also appears in Markman's
extended diagram preceding Lemma~8.3.4. The tangent bundle
is trivial, so its Todd correction is one.
If $\sigma_E$ is injective, the two observations imply
(S20.1). This fourfold criterion was also stated previously
as Theorem~A of the rStation paper~37, version~4, cited in the
source ledger; the argument here checks its conventions
against the supplied primary trace calculation.

We next compute the contraction rank. Put
$U=H^1(Y,\mathcal O_Y)$ and $T=H^0(Y,T_Y)$, and choose
dual bases so that a nondegenerate rational polarization
class is $\theta=\sum_{i=1}^4 x_i\wedge y_i$, with
$x_i\in H^0(\Omega_Y^1)$ and $y_i\in U$.
Let $s^2=-d$ and
\[
 w=\alpha+\beta
   =c_+e^{s\theta}+c_-e^{-s\theta},
 \qquad c_\pm=\frac{1\pm1/s}{2}.
\]
For
$\xi=(\eta,M,\nu)\in\wedge^2U\oplus(U\otimes T)
\oplus\wedge^2T=HT^2(Y)$, set
\[
 A=M\mathbin{\lrcorner}\theta\in\wedge^2U,
 \qquad B=\tfrac12\nu\mathbin{\lrcorner}\theta^2
             \in\wedge^2U,
 \qquad D=\eta-dB.
\]
The map $M\mapsto A$ is surjective and $\nu\mapsto B$ is
an isomorphism. The contraction convention is the one fixed
in (5.3) of the companion; reversing the bivector convention
merely changes the variable $B$ by its negative.
Exterior differentiation of the exponential gives
\[
 \xi\mathbin{\lrcorner}e^{\pm s\theta}
      =(D\pm sA)e^{\pm s\theta}.
\]
It follows that the first two Hodge blocks of
$\xi\mathbin{\lrcorner}w$ are
\[
 D+A,\qquad\theta\wedge(D-dA).
\]
The map $\theta\wedge:\wedge^2U\to H^3(\Omega_Y^1)$
is injective: if $\theta\wedge\gamma=0$, independence of
the $x_i$ gives $y_i\wedge\gamma=0$ for every $i$, hence
$\gamma=0$. The coefficient matrix on $(D,A)$ has
determinant $-(d+1)\ne0$. These two blocks therefore already
have rank twelve. Every block factors through the same two
six-dimensional spaces, so the full contraction rank is
exactly twelve.

The same rank holds for $e^{b\theta}w$, for any rational $b$.
Indeed the substitution
\[
 \eta' =\eta+bA+b^2B,\qquad A'=A+2bB
\]
is an invertible shear together with the unchanged variable
$B$, and reduces the two exponential slopes $b\pm s$ to
the preceding calculation after the invertible multiplication
by $e^{b\theta}$ on the target. This observation applies on
the quotient with the rational polarization whose pullback
is the original principal class.

It remains to verify the actual object. Write
$D_0=\bar\nu_*\mathcal O_{\widetilde Z/G}$.
It is a pure sheaf of dimension two: it injects into the
finite pushforward of the structure sheaves of the smooth
normalization components. The map $\mathcal O_Y\to D_0$
has kernel $I=\mathcal I_{\bar Z_{\mathrm{surf}}}$ and
zero-dimensional cokernel $Q$. Thus $\bar E$ has cohomology
$I$ in degree zero and $Q$ in degree one.
Since $I$ is torsion free,
$\operatorname{Hom}(Q,I)=0$. The truncation triangles show
that this implies
$\operatorname{Ext}^{j}(\bar E,\bar E)=0$ for $j<0$.
They also give
$\operatorname{Hom}(\bar E,D_0[-1])=0$, because
$\operatorname{Hom}(Q,D_0)=0$ by purity.
Codimension two of $\operatorname{Supp}D_0$ on the smooth
fourfold gives
$\operatorname{Hom}(D_0,\mathcal O_Y)
=\operatorname{Ext}^1(D_0,\mathcal O_Y)=0$.
Applying $\operatorname{Hom}(-,\mathcal O_Y)$ to
$\bar E\to\mathcal O_Y\to D_0$ therefore yields
$\operatorname{Hom}(\bar E,\mathcal O_Y)=\mathbf C$.
Applying $\operatorname{Hom}(\bar E,-)$ and using the other
vanishing embeds $\operatorname{End}(\bar E)$ into this
one-dimensional space. The identity is nonzero, so
$\bar E$ is simple.

Normalizing a transverse meeting point of two smooth
surfaces adds one skyscraper class to the structure sheaf.
The chosen four points for each of the $\binom\ell2$ pairs
therefore add $4\binom\ell2=2d\ell$ point classes upstairs.
This is exactly the number of isolated points in S-19.
Consequently the partial-normalization complex has the same
K-character as that first ideal seed, as asserted in
Markman Example~8.2.3. The quotient has
$\chi(\bar E,\bar E)=8d$, virtual rank one, and character
$e^{-\theta}w$ in the rational quotient polarization.
The preceding linear calculation gives
$\operatorname{rank}c_{\bar E}^2=12$.
The unit--trace composite is again rank one, so it injects
$H^1(\mathcal O_Y)$ into
$\operatorname{Ext}^1(\bar E,\bar E)$.
The negative-Ext vanishing, Serre duality, and simplicity
now give
\[
 \operatorname{ext}^2(\bar E,\bar E)
 =8d-2+2\operatorname{ext}^1(\bar E,\bar E)
 \ge8d+6>12.
\]
This contradicts (S20.1) if semiregularity is assumed.

For the quantitative assertion, put
$m=\operatorname{ext}^2(\bar E,\bar E)$ and
$a=\operatorname{rank}\sigma_{\bar E}
=\operatorname{rank}\operatorname{ob}_{\bar E}$.
The rank inequality for the composition gives
$12\ge2a-m$. Hence
\[
 \dim\ker\sigma_{\bar E}=m-a
 \ge\frac{m-12}{2}\ge4d-3.
\]
The argument proving simplicity also applies upstairs.
All linearizations therefore differ by characters of $G$,
and their descents by torsion line bundles, as in S-19.
Tensor equivalence preserves the nonzero full kernel,
proving the weak-equivariant assertion for this action.
\end{proof}

For both S-19 and S-20, taking an external product with a
nonzero perfect object preserves the displayed kernel:
the functor sends $u$ to $u\boxtimes\operatorname{id}$,
which is nonzero by K\"unneth, and the full trace is functorial.
The same is true after a derived equivalence.
Thus the canonical product and Orlov-transform route does not
repair these particular factors. This statement does not
exclude a different object or a different finite target action.

\begin{proposition}[S-21: All finite target actions obey the fourfold secant obstruction]
\label{prop:semireg-all-finite-target-actions}
Let $A$ be a complex abelian fourfold, $\mathcal C=
\operatorname{Perf}(A)$, and let a finite subgroup
$G\subset A\times\widehat A$ have a specified coherent action
on $\mathcal C$ by the corresponding translations and
tensor products with degree-zero line bundles. In particular,
no nonidentity group element represents the pair $(0,\mathcal O_A)$.
For a lift $\widetilde E\in\mathcal C^G$ of a perfect object
$E\in\mathcal C$, suppose
$\operatorname{Ext}^{j}_{\mathcal C^G}(\widetilde E,
\widetilde E)=0$ for $j<0$. Suppose further that
\[
 e^B\operatorname{ch}(E)=a\alpha+b\beta,
 \quad
 \alpha=\frac{e^{s\theta}+e^{-s\theta}}2,
 \quad
 \beta=\frac{e^{s\theta}-e^{-s\theta}}{2s},
 \quad s^2=-d,
 \tag{S21.1}
\]
where $d>0$, $\theta$ is a rational ample class,
$B$ is a rational algebraic divisor class, and
$a,b\in\mathbf Q$ are not both zero.
If $\widetilde E$ is semiregular in $\mathcal C^G$, then
\[
 \operatorname{ext}^2_{\mathcal C^G}(\widetilde E,
              \widetilde E)=12,
 \qquad
 \operatorname{ext}^0_{\mathcal C^G}(\widetilde E,
              \widetilde E)
 \ge 2+\frac{\chi_A(E,E)}{2|G|}>2,
 \tag{S21.2}
\]
where
\[
 \chi_A(E,E)=8d(da^2+b^2)\frac{\int_A\theta^4}{24}>0.
 \tag{S21.3}
\]
Consequently no such lift that is simple in the equivariant
category is semiregular. In particular, an underlying simple
perfect object with (S21.1) and no negative self-Ext admits no
weakly semiregular lift for any such finite target action.
The conclusion allows arbitrary mixtures of translations
and degree-zero tensor operations; it does not require a
Fourier--Mukai model on which the whole group acts freely
by translations.
\end{proposition}
\begin{proof}
Write $\phi_g$ for the functor associated to $g=(x,P)$.
The invariant category is smooth and proper in characteristic
zero; see Perry 2026, \S5, opening paragraph, p.~12,
and its cited Bayer--Perry Proposition~3.15.
First we verify that all nonidentity coefficient sectors vanish:
\[
 HH^*(\mathcal C,\phi_g)=0\qquad(g\ne1).
 \tag{S21.4}
\]
The coefficient space is the derived Hom from the diagonal
kernel to the kernel on the graph of translation by $x$,
with its line-bundle coefficient. This is the kernel
description in Perry, \emph{Hochschild cohomology and group
actions}, arXiv:1807.09268v2, Remark~4.1 and the proof of
Corollary~4.8, pp.~9 and~12. If $x\ne0$, the graph is
disjoint from the diagonal, so the derived Hom is zero.
If $x=0$ and $g\ne1$, then $P\ne\mathcal O_A$ and
the sheafified HKR decomposition identifies the coefficient
groups with sums of
$H^q(A,\wedge^pT_A\otimes P)$.
They vanish, because $T_A$ is trivial and a nontrivial
degree-zero line bundle on an abelian variety has no
cohomology. The latter also follows from the transform of
$\mathcal O_A$ in Mukai Theorem~2.2, p.~156, by derived
base change away from the origin.

Translations and degree-zero tensor operations act trivially
on $HH^*(A)$ and $HH_*(A)$: under HKR the former fix invariant
forms and vector fields, and the latter have zero Atiyah
class and Chern character one. Perry's Theorem~4.4 and
Remark~4.5, pp.~10--12, now give
\[
 HH^*(\mathcal C^G)\cong HH^*(\mathcal C).
 \tag{S21.5}
\]
The identity-sector inclusion is compatible with the
Hochschild algebra and with evaluation on equivariant objects.
This conclusion uses the specified coherent action, not
merely an unlinearized list of isomorphic functors.

The invariant category has Serre functor $[4]$.
To check this directly, the ordinary Serre trace pairing on
$\operatorname{Perf}(A)$ is invariant under translations
and degree-zero tensor products: translations preserve the
chosen invariant volume form, and tensoring preserves
perfect trace. Conjugation by the linearizations also
preserves trace. Taking $G$-invariants in the Serre duality
isomorphism remains a perfect duality, since $|G|$ is
invertible. This is exactly the invariant-duality argument
in the proof of Perry Lemma~6.5, equations (6.1)--(6.3),
p.~17; the proof there treats the general finite-group
Serre descent before its stated cyclic specialization.

Perry Definition~6.1 and Remark~6.2, p.~16, identify
Hochschild homology with Serre-coefficient Hochschild
cohomology. Converting his cohomological grading to the
usual homological convention, a category with Serre functor
$[4]$ has $HH_j\cong HH^{4-j}$ as vector spaces.
Thus (S21.5) gives equal finite dimensions for
$HH_j(\mathcal C^G)$ and $HH_j(\mathcal C)$, in particular
for $j=0,-2$.

Let $F:\mathcal C^G\to\mathcal C$ be forgetting and
$I:\mathcal C\to\mathcal C^G$ induction. These functors
are mutually left and right adjoint by Perry Lemma~3.3,
p.~7, and $FI\cong\bigoplus_{g\in G}\phi_g$.
Functoriality and additivity of Hochschild homology imply
\[
 F_*I_*=|G|\operatorname{id}_{HH_*(\mathcal C)}.
\]
One may read this additivity from Shklyarov, arXiv:0710.1937v3,
Proposition~3.2 and Theorem~3.4, pp.~14--16, by applying
the latter to a sum of functor kernels.
Equal dimensions then show that $F_*$ is an isomorphism
in degrees zero and minus two, and that
\[
 I_*F_*=|G|\operatorname{id}_{HH_j(\mathcal C^G)}
       \qquad(j=0,-2).
 \tag{S21.6}
\]
In particular, forgetting cannot hide a semiregularity
kernel in a nonidentity Hochschild sector.

We also need the exact Euler scaling, not an estimate by
averaging dimensions. Hochschild Chern characters commute
with functors, so (S21.6) implies
$\operatorname{ch}_{HH}(IF\widetilde E)
=|G|\operatorname{ch}_{HH}(\widetilde E)$.
The categorical Riemann--Roch formula makes the Euler
pairing linear in these characters: its precise source is
Shklyarov Theorem~3.5, p.~17, and its Hom formulation
equations (3.7) and (3.9), pp.~18--19.
The formula applies here via a proper DG-algebra model:
induction of a generator of $\operatorname{Perf}(A)$ is
a generator of $\mathcal C^G$, by adjunction and the
averaging retraction; its endomorphism algebra is proper.
Using the right adjunction $F\dashv I$, we obtain
\[
 \chi_A(E,E)
 =\chi_{\mathcal C^G}(\widetilde E,IF\widetilde E)
 =|G|\chi_{\mathcal C^G}(\widetilde E,\widetilde E).
 \tag{S21.7}
\]

Write $V_i=\operatorname{Ext}^i_A(E,E)^G$, which is the
degree-$i$ self-Ext in $\mathcal C^G$.
Perry 2026 Remarks~2.3 and~2.5, p.~5, give the commuting
semiregularity diagram for forgetting. Since $F_*$ in
degree minus two is an isomorphism, semiregularity of the
lift is equivalent to injectivity of
$\sigma_E|_{V_2}$.
The ordinary evaluation maps
$\operatorname{ob}_E^i:HH^i(A)\to\operatorname{Ext}^i(E,E)$
have image in $V_i$, because $G$ acts trivially on
$HH^i(A)$ and evaluation is natural.
The fourfold adjunction of S-20 restricts to the perfect
invariant Serre pairing on $V_2$. Consequently, if the
lift is semiregular, evaluation in degree two is onto
$V_2$ and
\[
 \dim V_2=\operatorname{rank}c_E^2.
 \tag{S21.8}
\]
Similarly $\dim V_1\ge\operatorname{rank}c_E^1$.
For clarity, the latter needs only the degree-one trace
identity. For $\gamma:\mathcal O_\Delta\to
\mathcal O_\Delta[1]$, naturality of the same adjunction
unit as in 032 equation (5.23) gives
\[
 (1_E\otimes L\Delta^*\gamma)AH_E
     =AH_E[1]\chi_E(\gamma).
\]
Take perfect trace with coefficients and apply HKR as in
(5.20)--(5.24). This gives
$\sigma_E^1\operatorname{ob}_E^1=c_E^1$,
so the claimed rank bound follows in the invariant image.

We calculate both contraction ranks explicitly. A
$B$-field factor $e^{-B}$ does not change them:
conjugation by $e^B$ replaces the degree-one contraction
operator $\iota_v$ by $\iota_v+\iota_vB\wedge$ and fixes
wedge operators. Since $B$ has type $(1,1)$, this is an
invertible shear of $HT^1(A)=U\oplus T$ and hence of its
exterior algebra, while $e^B$ preserves each Hodge-difference
space on the target.
For $w=a\alpha+b\beta$, the degree-two calculation in S-20
has first two blocks
\[
 aD+bA,\qquad\theta\wedge(bD-adA).
\]
The determinant is $-(a^2d+b^2)\ne0$, giving
$\operatorname{rank}c_E^2=12$.
In degree one, put $L=v\mathbin{\lrcorner}\theta\in U$
for $(\eta,v)\in U\oplus T$.
The first two blocks are
\[
 a\eta+bL,\qquad\theta\wedge(b\eta-adL).
\]
The same determinant is nonzero, and wedge by $\theta$
is injective on $U$. Thus $\operatorname{rank}c_E^1=8$.

Finally, $B$ cancels from $\operatorname{ch}(E)^\vee
\operatorname{ch}(E)$. If
$c_\pm=(a\pm b/s)/2$, then
$c_+c_-=(a^2+b^2/d)/4$.
The top term of $w^\vee w$ is
$2c_+c_-(2s)^4\theta^4/24$, proving (S21.3).
Put $h=\dim V_0$. The assumed negative-Ext vanishing and
invariant Serre duality imply vanishing above degree four,
$\dim V_4=h$, and $\dim V_3=\dim V_1$.
Equations (S21.7)--(S21.8) therefore imply
\[
 12=\dim V_2
 =\frac{\chi_A(E,E)}{|G|}-2h+2\dim V_1
 \ge\frac{\chi_A(E,E)}{|G|}-2h+16.
\]
Rearranging gives (S21.2), including its contradiction
when $h=1$. The proof applies separately to every coherent
linearization; none was assumed canonical.
\end{proof}

The earlier fourfold criterion and its free-translation
version are recorded in rStation paper~37, version~4,
Theorems~A, B and $B_G$. S-21 adds the coefficient-sector
and induction argument for arbitrary effective finite
subgroups of $A\times\widehat A$, and permits simplicity
only after equivariantization. It excludes this ample
imaginary-quadratic secant character for equivariantly
simple seeds; it does not exclude arbitrary corrected
characters on the final abelian eightfold.
